

Using Theorem 12.9 about marginal distributions we get that 

\[
    f_{X_1} (x) = \int_{\R}^{} f_{X} (x,y) dy
\] 
and 
\[
    f_{X_2} (x) = \int_{\R}^{} f_{X} (x,y) dx
.\] 
Let's first compute the density of the marginal distribution of the random variable $ X_1$.
If $ x \in [0,1]$ then

\begin{align*}
    f_{X_1} (x) &= \int_{\R}^{} (x^2y + 1/12 y^{3} )dy \\
		&= 2x^2 + 1/3
\end{align*}
If $ x \notin [0,1]$ then $ f_{X} (x,y) = 0$ and we get that the integral over $ \R$ is zero meaning
that $ f_{X_{1} } (x) = 0$.
Combining these results we get that 

\[
    f_{X_1} (x) = \begin{cases}
	2x^2 + 1/3 &, \, x \in [0,1]\\ 
	0 &, \, \mbox{else} 
    \end{cases}
.\] 

Now let's find the density for the random variable $ X_2$.
If $ y \in [0,2]$ then

\begin{align*}
    f_{X_2} (y) &= \int_{\R}^{} x^2y + 1/12 y^{3} dx \\
		&= y [x^{3}/3]_{0}^{1} + 1/12 y^{3}[x]_{0}^{1} \\
		&= \frac{y}{3} (1+ \frac{y^2}{4} )
\end{align*}.
Again, if $ y \notin [0,2]$ then $ f_{X_2}(y) = 0$ since $ f_{X}(x,y) = 0$.
Combining these we get

\[
    f_{X_2} (y ) = \begin{cases}
	\frac{y}{12} (4 + y^2) &, \, y \in [0,2]\\
	0 &, \, y \notin [0,2]
    \end{cases}
.\] 
